\section{Introduction}
Automated electrocardiographic analysis is most useful when it can distinguish clinically relevant rhythm changes without confusing patient-specific morphology for disease. That requirement is difficult to satisfy in heartbeat classification because ambulatory ECG data combine substantial inter-patient variation with highly unequal class prevalence: normal beats dominate most recordings, whereas supraventricular ectopic and fusion beats may be represented by only a small fraction of the available examples. Under these conditions, a high aggregate score can coexist with poor recognition of the rhythms for which automated review is intended to add value\cite{chawla2002smote,he2009learning,saito2015precision}. The problem is therefore not simply to increase classification accuracy, but to determine whether a learned representation remains informative across patients and whether its errors are clinically concentrated in particular rhythm classes.

Evaluation protocol is central to that question. ECG beats sampled from the same patient share morphology, acquisition characteristics, and rhythm context; consequently, beat-level random partitioning can answer a different and generally easier question than evaluation on records from patients excluded from model development. The literature demonstrates the potential of convolutional, recurrent, residual, and patient-adapted models for arrhythmia analysis\cite{kiranyaz2015personalized,luz2016ecg,hannun2019cardiologist,acharya2017deep}, but direct comparison of headline metrics is complicated by differences in databases, label definitions, segment duration, and partitioning strategy\cite{ismailfawaz2019deeplearning}. For a reproducible benchmark, the split, preprocessing path, comparator implementations, and class-specific outcomes must therefore be reported with the same prominence as the final accuracy.

From a representation perspective, heartbeat morphology contains both local time--frequency structure and context extending across adjacent beats. Wavelet transforms provide a multiresolution description of transient ECG morphology\cite{addison2005wavelet,torrence1998practical,mallat1989wavelet}, while convolutional layers offer a strong local inductive bias and self-attention can integrate information over a broader sequence. Conformer encoders combine these operations through convolution, multi-head self-attention, and feed-forward modules\cite{gulati2020conformer}. This architecture is therefore a plausible candidate for Morlet-scalogram ECG sequences, but its value should be established through controlled comparison rather than architectural novelty alone. Equally important, any gain should be examined alongside minority-class failure, because moderate ranking discrimination can still yield an unusable operating point when prevalence is low.

In this study, we evaluate a Morlet CNN--Conformer pipeline on the MIT-BIH Arrhythmia Database using a fixed record-level development/test partition that preserves subject separation at final evaluation\cite{goldberger2000physiobank,moody2001impact}. The contribution is threefold. First, we provide an auditable preprocessing, training, and evaluation pipeline in which the held-out cohort is isolated from model fitting and hyperparameter selection. Second, we compare the proposed configuration with attention-only, CNN--LSTM, and feature-engineered alternatives under the same held-out split. Third, we make class imbalance and failure analysis part of the primary result by reporting class support, precision, recall, F1, AUC, dominant confusion pathways, threshold sweeps, and internal post-hoc calibration diagnostics. This framing is intended to establish what the architecture improves, what remains unresolved for rare rhythms, and which evidence is still required before considering clinician-in-the-loop evaluation.